\begin{proposition}[Explicit height enclosures]\hypertarget{tag-000Z}{}\label{ari:C.6.5}
\StatementRecord{ari:C.6.5}{proved}{000Z}{\SeriesRef{ari:C.6.3}, \SeriesRef{ari:C.5.1}}

For each of the eleven rational CM parameters in \cref{ari:C.6.3},
the Fourier and product algorithms below give rational interval
enclosures of $\log N(P)$ and $\widehat h_O(P)$ of radius less
than $10^{-236}$.
\end{proposition}


\begin{proof}
\begin{enumerate}[label=\arabic*.,leftmargin=*]
\item The CM discriminants, in the order of \ref{ari:C.6.3}, are
\[
 -8,-12,-16,-28,-7,-11,-19,-27,-43,-67,-163.
\]
Enumeration of primitive reduced forms
$(a,b,c)$, $|b|\le a\le c$, $b^2-4ac=D$, with the usual
nonnegative boundary convention, gives exactly one form of
each discriminant, always with $a=1$.
The range is finite and proved by $a\le\sqrt{|D|/3}$.
Thus the CM theorem used in \ref{ari:C.5.1} makes the corresponding
singular modulus an integer. The Fourier enclosures below
place it within $1/4$ of the specified $j$, proving its
identity. 

\item Write $\tau=iT$ for the first four entries and
$\tau=1/2+iT$ for the other seven, with
\[
 (T_1,T_2,T_3,T_4)=(\sqrt2,\sqrt3,2,\sqrt7),\qquad
 T_k=\sqrt{|D_k|}/2\quad(k\ge5).
\]
Put $Q=e^{-2\pi T}$, $q=Q$ or $-Q$ accordingly,
$M=350$, and $Y=2/5$. On $0\le y\le Y$,
\begin{equation}\hypertarget{tag-003C}{}\label{ari:eq:ch-wp-bracket}
 \frac{\wp(1/2-iy)}{\pi^2}
 =\frac1{\cosh^2\pi y}-\frac13+
 8\sum_{n\ge1}\frac{nq^n}{1-q^n}
                 \{1-(-1)^n\cosh(2\pi ny)\}.
\end{equation}
For $a=Qe^{2\pi Y}$, the absolute tail after $M$ in
$\wp$ is at most
\[
 \frac{16\pi^2}{1-Q}
 \frac{a^{M+1}\{M+1-Ma\}}{(1-a)^2}<10^{-875}.
\]
The proof is $1+\cosh(2\pi ny)\le2e^{2\pi nY}$ and
$\sum_{n>M}na^n=a^{M+1}(M+1-Ma)/(1-a)^2$.
A numerical zero approximation selects a rational centre $y_0$. At the two rational endpoints
$y_0\pm10^{-240}$ the truncated expression plus this
infinite tail is respectively strictly positive and strictly
negative. These strict endpoint signs are evaluated with outward error bounds.
Continuity produces a zero within each bracket. Every zero
of $\wp$ is one of the pair $\pm P$, and $N$ agrees on
that pair, so no unproved choice of root is needed.

\item The product formula on this bracket has positive real factors:
\begin{align*}
 \log N={}&\log(2\cosh\pi y)-\frac{\pi T}{6}
                      -\frac{\pi y^2}{T}\\
 &+\sum_{n\ge1}\{\log(1+q^ne^{2\pi y})
                         +\log(1+q^ne^{-2\pi y})\}.
\end{align*}
Since $|\log(1+x)|\le |x|/(1-|x|)$ for real $|x|<1$,
the omitted tail is bounded by
\[
 \frac{2e^{2\pi Y}Q^{M+1}}
 {(1-Q)(1-e^{2\pi Y}Q^{M+1})}<10^{-1260}.
\]
The finite sums, elementary functions and tails are evaluated by
outward interval arithmetic. The radius and sign comparisons are strict. The resulting radius is less than $10^{-236}$.
Adding the exact finite local logarithms of \ref{ari:C.6.3} and
subtracting this enclosure gives the claimed height bound.

\item The singular-modulus certificate uses
$E_4=1+240\sum n^3q^n/(1-q^n)$ and
$\Delta_0=q\prod(1-q^n)^{24}$. Their respective error
bounds, after truncation at $M$, are
\[
 \frac{240(M+1)^3Q^{M+1}(1+4Q+Q^2)}{(1-Q)^5},
 \qquad
 \frac{24Q^{M+1}}{(1-Q)(1-Q^{M+1})}
\]
for $E_4$ and for the logarithm of the product. The first
uses $(M+1+k)^3\le(M+1)^3(k+1)^3$ and the differentiated
geometric series; the second uses the same logarithm bound.
Thus $E_4^3/\Delta_0$ is enclosed by the stated finite sums and tails.
\end{enumerate}
\end{proof}